\chapter{Introduction} 

\label{cha:introduction}
%==============================================================================

%\blindtext

%\section{Motivation}
\label{sec:motivation}
%==============================================================================

\blindtext \cite{Loepmeier:2015TR} \label{example:reference}

%http://get-software.net/macros/latex/contrib/caption/subcaption.pdf
\begin{figure}	
	\centering
	\begin{subfigure}[t]{0.45\textwidth}
		\caption{Logo des Institut für Informatik}
		\includegraphics[width=1.0\textwidth]{graphics/cs_logo_nl_blue.pdf}
	\end{subfigure}
	\begin{subfigure}[t]{0.45\textwidth}
		\caption{Logo der Uni}
		\includegraphics[width=1.0\textwidth]{graphics/wwu-logo-neu.pdf}
	\end{subfigure}

	\begin{subfigure}[t]{0.45\textwidth}		
		\phantomsubcaption\label{fig:inner}
		\missingfigure[figwidth=\textwidth]{Logo des FB10}
	\end{subfigure}
	\begin{subfigure}[t]{0.45\textwidth}
		\phantomsubcaption
		\missingfigure[figwidth=\textwidth]{Logo der Numerik}		
	\end{subfigure}
\caption[Logos der Uni]{Logos der Uni. Dabei kann man auch wieder referenzieren wie auf \subref{fig:inner}. Dabei muss noch nicht einmal eine Unterüberschrift existieren. Dies ist besonders praktisch, falls die Buchstaben schon in den Bildern vorhanden sind. Details \href{http://get-software.net/macros/latex/contrib/caption/subcaption.pdf}{hier}
}
\label{fig:meth:bgsub:whole}
\end{figure}
%
%\Blindtext
%
%\section{Thesis Structure}
%\label{secStruktur}
%%==============================================================================
%
%\Blindtext
